% Included inside the book's existing thebibliography environment.
% Nam (2024) is already entered as nam2024 in util/stative-references.tex.
% Access and claim limits: util/stative-extension-source-audit.md.
\bibitem{song2018lexical}
Song, Jaemog. 2018.
Lexical Evidentiality in Korean.
\textit{Eoneohag} 82: 105--131.
The Linguistic Society of Korea; ISSN 1225-7494.
\href{https://doi.org/10.17290/jlsk.2018..82.105}{doi:10.17290/jlsk.2018..82.105}.

\bibitem{park2026perspective}
Park, Myung-Kwan. 2026.
Perspective Effects on Korean Subjects in Jussive / Commissive Endings,
Evidentials, and Psych Predicates.
\textit{Studies in Generative Grammar} 36(2): 165--196.
\href{https://doi.org/10.15860/sigg.36.2.202605.165}{doi:10.15860/sigg.36.2.202605.165}.

\bibitem{kimyoonshin2026}
Kim, Yoon-shin. 2026.
A Study on Event Structure of Psych-predicates: Focusing on Korean
Psych-adjectives and `-ehata' Type Psych-verbs.
\textit{Eoneohag} 105: 3--24.
The Linguistic Society of Korea; ISSN 1225-7494.
\href{https://doi.org/10.17290/jlsk.2026..105.3}{doi:10.17290/jlsk.2026..105.3}.

\bibitem{frellesvig2010}
Frellesvig, Bjarke. 2010.
\textit{A History of the Japanese Language}.
Cambridge: Cambridge University Press. Published online 3 May 2011.
\href{https://doi.org/10.1017/CBO9780511778322}{doi:10.1017/CBO9780511778322}.

\bibitem{wandl2023univerbation}
Wandl, Florian. 2023.
Univerbation and prosodic change: On the origin of the Slavic definite adjective
accentuation.
\textit{Indo-European Linguistics} 11(1): 241--287.
\href{https://doi.org/10.1163/22125892-bja10024}{doi:10.1163/22125892-bja10024}.
